% ----- Consignes exo 3 ----- %
\begin{td-exo}[Autour du \textsc{Max} \(3\)-\textsc{Dimensional Matching}]\,\\ % 3 
	On rappelle la formulation du problème comme suit:

	\begin{problem}{\textsc{Max} \(3\)-\textsc{Dimensional Matching} (MAXDM)}
    \Input{Un ensemble \(S \subseteq X \times Y \times Z\) de triplets.}
    \Question{Trouver un couplage de cardinalité maximum, i.e.\ un ensemble \(M \subseteq S\) où toute paire de triplets de \(M\) n'ont aucun élément en commun.}
\end{problem}

	et on propose l'algorithme d'approximation suivant:

	\begin{styledalgo}{Algorithme pour le Max 3D Matching}\label{alg:exam1_ex03_1}
    \Input{Un ensemble \(S \subseteq X \times Y \times Z\) de triplets.}
    \Output{Un couplage \(M \subseteq S\)}
    \State \(M \gets \varnothing\) % chktex 1
    \For{chaque triplet \((x,y,z) \in S\)}
        \If{\((x,y,z)\) ne contredit pas \(M\)}
            \State \(M \gets M \cup \{(x,y,z)\}\) % chktex 1
        \EndIf % chktex 1
    \EndFor % chktex 1
    \State \textbf{return} \(M\) % chktex 1
\end{styledalgo}


	\begin{enumerate}
		\item Donner la complexité de l'algorithme 1 % TODO: add direct reference to algo
		\item Montrer que l'algorithme 1 est \(3\)-approché. % TODO: add direct reference to algo
		\item Trouver une instance pour laquelle la borne de trois est atteinte.
		
		\textbf{Aide:} Considérer un ensemble de \(4h\) triplets avec \(h \in \mathbb{N}\).
		L'algorithme prendra que \(h\) triplets tandis que la solution optimale en prendra \(3h\).
	\end{enumerate}
\end{td-exo}

% ----- Solutions exo 3 ----- %
\iftoggle{showsolutions}{ 
	\begin{td-sol}[]\ % 3
		\begin{enumerate}
			\item L'algorithme a une complexité de \(O(n^2)\).

			\item L'algorithme ne peut pas etre mieux que 3 approché, car on a un exemple de cas 3-approché. 
			De plus, a chaque étape on dégage au pire 3 triplets de la solution optimale, alors on a forcément une approximation de facteur 3.

			\item Il suffit de considérer 4 triangles dont 3 qui connectent à un sommet du triangle restant.
		\end{enumerate}
	\end{td-sol}
}{}
